\subsection{達成成果のサマリ}
\label{sec:achievements_summary}
% ═══════════════════════════════════════════════════════════════════════════════

本稿の三本柱設計は，単体検証（第\ref{sec:numerical_integrity}章），統合検証
（第\ref{sec:verification}章），および物理ベンチマーク
（第\ref{sec:validation}章）で段階的に評価した．
第\ref{sec:validation}章では，毛管波，振動液滴，気泡上昇，
Rayleigh--Taylor 不安定性を同一標準構成上に置き，
復元振動，閉界面保存量ゲート，浮力上昇，不安定成長という
異なる物理応答を読み分けた．

% ─────────────────────────────────────────────────────────────────────────────
\subsubsection{単体検証 U1--U12 の総括}
% ─────────────────────────────────────────────────────────────────────────────

第\ref{sec:numerical_integrity}章の単体検証は 11 スイート 26 サブテストから構成され，
集計（表~\ref{tab:U_summary}）は \textbf{23\,\checkmark + 2\,$\triangle$ + 1\,$\bullet$，$\times$ ゼロ}である．

\begin{description}
  \item[U1（CCD 演算子群；§\ref{sec:U1_ccd_suite}）：]
        CCD 内点 $\Ord{h^{6.0}}$，DCCD ナイキスト除去 $|H(\pi)|=0$（機械精度），
        FCCD 面値・面勾配 $\Ord{h^{4.00/4.00}}$，UCCD6 スーパー収束 $\Ord{h^{7.00}}$ を実証．
  \item[U2（CCD-Poisson PPE 境界条件；§\ref{sec:U2_ccd_poisson_ppe_bc}）：]
        標準設定で $\Ord{h^{5.72}}$，欠陥補正 DC k=3 で $\Ord{h^{6.90}}$ を達成．
  \item[U3（非一様格子上の空間離散；§\ref{sec:U3_nonuniform_suite}）：]
        非一様 CCD $\Ord{h^{5.98}}$，FCCD $\Ord{h^{4.00/4.09}}$，
        GCL 機械零（$2.13\times 10^{-13}$），Ridge--Eikonal D1--D4 機械零を確認．
  \item[U4（Ridge--Eikonal 再初期化；§\ref{sec:U4_ridge_eikonal_reinit}）：]
        Godunov Eikonal + DGR で $\varepsilon_\mathrm{eff}/\varepsilon_* = 2.0\to 1.0000033$ の
        SDF 距離関数品質回復を実証．
  \item[U5（Heaviside / $\delta_\varepsilon$ モーメント；§\ref{sec:U5_heaviside_delta}）：]
        $\int H_\varepsilon\,dx$ と $\int \delta_\varepsilon\,dx$ のモーメント残差が機械精度 floor に到達．
  \item[U6（分相 PPE + HFE + DC；§\ref{sec:U6_split_ppe_dc_hfe}）：]
        分相 PPE + DC 残差 $7\times 10^{-6}$，HFE Hermite 場延長
        1D $\Ord{h^{5.91}}$，2D effective slope max $11.35$ med $8.72$ を確認
        （形式次数主張ではなく丸め床を含む診断値）．
  \item[U7（BF 静止液滴；§\ref{sec:U7_bf_static_droplet}）：]
        Laplace 圧力誤差 $0.61\%$ @ $N=128$，
        面 $\mu$ 補間（harmonic / arithmetic / log）の収束 slope $1.91/1.81/1.91$．
  \item[U8（時間積分；§\ref{sec:U8_time_integration}）：]
        TVD-RK3 $\Ord{\Delta t^{3.01}}$，AB2+IPC $\Ord{\Delta t^{2.04}}$，
        CN ADI $\Ord{\Delta t^{2.00}}$．3 層粘性 split は条件付き $\triangle$．
  \item[U9（DCCD-on-pressure 禁止の反証；§\ref{sec:U9_dccd_pressure_prohibition}）：]
        圧力に対する DCCD 適用は固有モード比 $10^{5}$--$3\times 10^{8}$ で
        圧力場を破壊するため禁止が成立（$\bullet$，研究上の反例として記録）．
  \item[U10（保存形共通フラックス台帳；§\ref{sec:U10_common_flux_ledger}）：]
        closed common-flux candidate は $\max|\Delta M_\rho|=1.137\times10^{-13}$，
        $\max|\Delta P_i|=2.842\times10^{-14}$，アフィン密度誤差 $0$ を達成し，
        非アフィン密度と $\psi$-only ledger の負対照を $2/2$ reject した．
  \item[U11（制約付き面速度状態空間；§\ref{sec:U11_constrained_face_state_space}）：]
        wall trace $\le2.581\times10^{-12}$，冪等性 $4.851\times10^{-12}$，
        Green residual $\le1.294\times10^{-16}$，階数ゲート wall $19/19$，
        periodic-wall quotient $25/25$ を確認した．
  \item[U12（アクティブ幾何毛管分解ゲート；§\ref{sec:U12_ao_capillary_split_gate}）：]
        全圧力像は毛管波の釣合い駆動を $0$ にする反例であり，
        成分体積 Hodge 診断は N32/N64 の非静的釣合いノルム
        $2.1176/2.3055$ を検出した．
        $R_p$ 未確定の分解，圧力履歴座標を持たない分解，
        および移動格子の投影済み面共鎖移送を欠く分解は
        標準経路へ入れない境界として整理した．
\end{description}

% ─────────────────────────────────────────────────────────────────────────────
\subsubsection{統合検証 V1--V11 の総括}
% ─────────────────────────────────────────────────────────────────────────────

第\ref{sec:verification}章の統合検証は 11 実験から構成され，
集計（表~\ref{tab:V_summary}・表~\ref{tab:v_accuracy_summary}）は $\times$ ゼロを維持している．
V1-a/V1-b/V2/V4/V5 は判定基準を実装 identity に合わせて見直し，
それぞれ FVM PPE 結合 stack の $O(h^2)$ trend，standard PPE $O(\Delta t)$，
Kovasznay wall residual $O(h^4)$，fixed-wall offset residual $O(\|\bm U\|)$，
CCD 寄生流れ絶対値 $u_\infty^{\mathrm{end}}$ の $\sigma/\mathrm{Re}_c$ スケール基準
に対する合格へ再分類した．
主張支持・条件付き採用実験（V1--V9）の headline と，
条件付き CLS 移流ベースライン（V10-a/b）は以下：

\begin{description}
  \item[V1（TGV；§\ref{sec:energy_conservation}）：]
        空間 trend は $E_{\mathrm{rel}}=1.583\times10^{-7}$ @ $N=64$，
        V1-b は standard PPE 結合の時間 slope $1.00$．
        AB2 単体 2 次ではなく縮約 projection stack の $O(\Delta t)$ 診断として合格．
  \item[V2（Kovasznay residual；§\ref{sec:kovasznay}）：]
        wall one-sided 境界閉包を含む統合 residual で slope $3.95$．
        CCD 内点 6 次ではなく Kovasznay wall residual $O(h^4)$ 診断として合格．
  \item[V3（静止液滴 long-term）：]
        $|\Delta p_\mathrm{rel}| = 0.97\%$ @ $N=128$．Laplace 圧力跳びの長期維持を確認．
        標準 pressure-jump/Riesz 静止ゲートでは
        $N=32,T=10$ で $K_{\max}=1.718\times10^{-6}$，
        $|\Delta V|/V_0\le2.919\times10^{-15}$，$D=0$ を確認した．
  \item[V4（fixed-wall uniform-offset；§\ref{sec:galilean_offset}）：]
        $\|\Delta u\|_\infty^{\mathrm{end}}=2.07\times10^{-1}$，
        $\|\Delta u\|_\infty^{\max}/\|\bm U\|=2.63<3$．
        厳密ゼロ残差ではなく fixed-wall uniform-offset residual として合格．
  \item[V5（寄生流れ multistep；§\ref{sec:spurious_multistep}）：]
        CCD $u_\infty^{\mathrm{end}}=1.93\!\times\!10^{-2}$ @ $\rho_r=1$, $N=32$（FD 比 $1/5.44$ 補助）．
        CCD は寄生流れ抑制に有効．
  \item[V6（密度比 sweep；§\ref{sec:density_ratio_sweep}）：]
        $\rho_r\in\{2,10,100,833\}$ の全 8 ケースでブローアップなし，
        $u_\infty^{\mathrm{end}}\le1.3\times10^{-9}$，
        $|\Delta V_\psi|/V_{\psi,0}\le1.2\times10^{-16}$．
        分相 PPE + DC + HFE + capillary range projection が高密度比でも機能することを実証．
  \item[V8（非一様格子 NS 静止液滴；§\ref{sec:nonuniform_ns_static}）：]
        $|\Delta p_\mathrm{rel}| = 2.86\%$ @ $\alpha=2$, $N=96$．固定非一様格子上の静止液滴として，
        非一様 D1--D4 と FCCD の整合を確認．
  \item[V9（local-$\varepsilon$ switch；§\ref{sec:local_eps_validation}）：]
        range-projected pressure-jump stack 上で B/C=$1.00$ @ $N=32$，
        $u_\infty^{\max}\le9.82\times10^{-10}$，
        $|\Delta V|/V_0\le4.73\times10^{-16}$．空間可変 $\varepsilon(\bx)$ は
        pressure-jump/HFE 短時間 probe に限定し，CSF 経路と汎用長時間移動界面からは分離．
  \item[V10-a（Zalesak; §\ref{sec:zalesak_cls_advection}）：]
        $\alpha=1$ の一様格子で質量保存軸 mass drift $4.70\!\times\!10^{-4}\%$ @ $N=128$ \checkmark (Type-B)，
        形状軸 centroid err $4.911\!\times\!10^{-3}$ ($\triangle$ Type-D, fixed-grid 位相/threshold 誤差と slot/h$=6.4$ 律速)．
        詳細は表~\ref{tab:V10_zalesak}．
  \item[V10-b（単一渦反転; §\ref{sec:single_vortex_reversal}）：]
        $\alpha=1, N=128, T=8$ で質量保存軸 mass drift $0.0428\%$ \checkmark (Type-B)，
        形状復元軸 $L^1$ reversal error $2.248\times 10^{-2}$
        ($\triangle$ Type-D, grid-scale folded filament の Eulerian 固定格子下限)．
        いずれも Olsson--Kreiss 質量修正 10 step ごと適用後．
  \item[V11（アクティブ幾何毛管分解の統合ゲート；§\ref{sec:ao_fast_capillary_gate}）：]
        U12 の反例を統合検証文脈へ接続し，
        §\ref{sec:ao_fast_state_space} の
        $R_p$ 分解，幾何互換，Hodge 逆適用済み面補間，移動格子面共鎖，圧力履歴座標，
        DC 収束判定，近似誤差境界を標準物理経路の採用境界として固定した．
\end{description}

条件付き $\triangle$ は V7 IMEX-BDF2 effective slope $1.59$ coupled-stack 律速
（Type-D, capillary pressure-jump/projection 界面帯に律速され BDF2 単体 $2.0$ に届かない），
V10-a Zalesak $\alpha=1$ 形状軸 (centroid err $4.911\!\times\!10^{-3}$，fixed-grid 位相/threshold 誤差と slot/h=6.4 律速；
質量保存軸は \checkmark)，
V10-b 単一渦 reversal $L^1=2.248\times10^{-2}$ に限定する．
V1-a/V1-b/V2/V4/V5 は本章の revised criterion により合格へ再分類済みであり
（V1-a は単点 $E_{\mathrm{rel}}$ ではなく空間 slope $2.00$ on FVM PPE $O(h^2)$ trend，
V5 は CCD/FD 比ではなく CCD 寄生流れ絶対値 $u_\infty^{\mathrm{end}}$ の $\sigma/\mathrm{Re}_c$ スケール基準），
今後 IPC/rotational projection，6 次境界閉包，moving-wall/ALE Galilean test へ拡張する際の
定量的ベースラインを与えている．
残る V7/V10 は格子増強だけでなく，capillary pressure-jump/projection，norm 選択，
参照解品質，再初期化間隔，非可逆な形状復元を分離する
追加検証ゲートとして残す．

% ─────────────────────────────────────────────────────────────────────────────
\subsubsection{物理ベンチマーク §\ref{sec:validation} の総括}
% ─────────────────────────────────────────────────────────────────────────────

第\ref{sec:validation}章の物理ベンチマーク 4 種
（毛管波・振動液滴・気泡上昇・Rayleigh--Taylor）は，
古典理論・既往 DNS との比較対象，実行条件，出力診断，判定軸を固定している．
表~\ref{tab:ch14_summary} の通り，4 ケースはいずれも同じ
active-geometry capillary decomposition / FCCD PPE / UCCD6 / pressure-coordinate history / IMEX--BDF2 構成を用い，
物理機構だけを変えて読むベンチマークである．
ただし振動液滴は，1 周期再実験で有限には完走したものの
鋭界面体積保存ゲートを満たしておらず，閉界面の保存量契約を
未解決課題として切り出す未受理ケースである．

\begin{description}
  \item[毛管波（§\ref{sec:val_capillary}）：]
        平坦界面の小振幅復元力を用い，
        bundle Young--Laplace work と pressure-component Hodge reaction が
        restoring 向きに PPE/corrector へ入ることを確認した．
  \item[気泡上昇（§\ref{sec:val_rising_bubble}）：]
        Hysing ら~\cite{Hysing2009} の単一気泡幾何を参照しつつ，
        水--空気級密度比・非一様格子・pressure-jump・浮力残差分解を同時に有効化する
        閉包ベンチマークとして位置付ける．
        $y_c(t)$，平均上昇速度，変形量，体積保存を同時に読み，
        hydrostatic-balanced buoyancy が投影と同じ面速度空間で閉じることを確認した．
  \item[振動液滴（§\ref{sec:val_oscillating_droplet}）：]
        静止液滴は第\ref{sec:verification}章の BF 検証に置き，
        本章では水--空気楕円液滴の Rayleigh--Lamb 型振動を扱う．
        実行条件は $32^2$ periodic，$\alpha=2$，
        $(a,b)=(5.5,4.5)\,\mathrm{mm}$，$D_0=0.10$，
        $T=0.105460663082\,\mathrm{s}$ とし，
        \texttt{signed\_deformation} の位相，粘性減衰包絡，
        \texttt{kinetic\_energy} の交換，体積保存を読む．
        再実験は 3452 samples で完走したが，
        最終 \texttt{volume\_conservation}=3.55e-1，
        P1 sharp 面積は $7.7395\times10^{-5}$ から
        $9.7435\times10^{-5}$ へ +25.9\% drift したため，
        検証合格ではなく閉界面保存量契約の残課題と判定した．
  \item[Rayleigh--Taylor 不安定（§\ref{sec:val_rayleigh_taylor}）：]
        $\mathrm{At}=0.0256$, $k=4\pi$ の線形成長率
        $\gamma \approx 1.78\,\mathrm{s^{-1}}$（Chandrasekhar~\cite{Chandrasekhar1961}
        古典理論）を基準に，$0 < t < 3/\gamma$ の対数線形窓，
        $T=7$ の非線形 mushroom 形成，体積保存，運動エネルギー推移を確認した．
\end{description}

これら 4 ベンチマークは，三本柱設計（CCD 高次群／分相 PPE+HFE+DC／BF+FCCD 面ジェット）の
物理的有効性を確認するための最終段である．
部品検証で確定した離散契約は，安定振動，浮力上昇，
不安定成長という異なる現象に対して一貫した応答を与える一方で，
閉界面振動では同じ cochain 上の sharp 体積保存がまだ破れていることを示した．

% ═══════════════════════════════════════════════════════════════════════════════
